% !TEX program = xelatex
% =====================================================================
%  通用報告樣板 (XeLaTeX + 中文 TW-Kai,含粗體)
%  編譯:xelatex -> biber -> xelatex -> xelatex   (或 latexmk -xelatex)
%  需求:TW-Kai(全字庫正楷體)字型。Overleaf 內建;本機可從
%        https://www.cns11643.gov.tw/opendata/Fonts_Kai.zip 取得後
%        放入 ~/.local/share/fonts/ 並執行 fc-cache。
% =====================================================================
\documentclass[12pt,a4paper]{article}

% ---- 版面 ----------------------------------------------------------
\usepackage[margin=1in]{geometry}
\usepackage{setspace}
\onehalfspacing
\usepackage{titling}
\setlength{\droptitle}{-1.5cm}
\usepackage{fancyhdr}
\setlength{\headheight}{15pt}
\pagestyle{fancy}\fancyhf{}
\renewcommand{\headrulewidth}{0.4pt}
\lhead{\leftmark}\rhead{\thepage}
\parindent=24pt

% ---- 數學與符號 ----------------------------------------------------
\usepackage{amsmath,amsthm,amssymb,amsbsy,bm}
\usepackage{siunitx}

% ---- 圖、表 --------------------------------------------------------
\usepackage{graphicx}
\graphicspath{{figures/}{reports/figures/}}
\usepackage{float}
\usepackage{subfig}
\usepackage{booktabs}
\usepackage{array}
\usepackage{multirow}
\usepackage{longtable}
\usepackage{enumitem}

% ---- 連結 ----------------------------------------------------------
\usepackage{xcolor}
\usepackage[unicode=true,pdfborder={0 0 0},bookmarksdepth=-1,
            colorlinks=true,linkcolor=blue!55!black,citecolor=blue!55!black,
            urlcolor=teal]{hyperref}

% ---- 參考文獻(biblatex + biber)----------------------------------
\usepackage[backend=biber,style=numeric,giveninits=true]{biblatex}
\addbibresource{references.bib}
% 範例書目(自含,讓樣板可獨立編譯;實際使用時刪掉此 filecontents、
% 改用你自己的 references.bib):
\begin{filecontents}[overwrite]{references.bib}
@article{Sample2024,
  author  = {Wang, M. and Lee, C.},
  title   = {An example reference for the template},
  journal = {Journal of Examples},
  year    = {2024}, volume = {1}, pages = {1--10}
}
\end{filecontents}

% ---- 中文字型:TW-Kai,並啟用「假粗體 / 假斜體」-------------------
%  TW-Kai 僅有 Regular 一個字重,故以 AutoFakeBold 合成粗體
%  (\textbf{中文} 才會變粗),AutoFakeSlant 合成斜體。
\usepackage{xeCJK}
\setCJKmainfont{TW-Kai}[AutoFakeBold=2.5, AutoFakeSlant=0.2]

% ---- 自訂指令 / 定理樣式 ------------------------------------------
\newcommand{\gauss}{\ \text{G}}
\newtheoremstyle{mystyle}{6pt}{15pt}{}{}{\bfseries}{.}{1em}{}
\theoremstyle{mystyle}
\newtheorem{thm}{定理}
\newtheorem{dfn}{定義}

% =====================================================================
\title{\bfseries 報告標題(中文粗體)\\[4pt]
       \large 副標題}
\author{作者一、作者二\\ \normalsize 國立臺灣大學 物理學系}
\date{實驗日期:2026 年 5 月 18 日}

\begin{document}
\maketitle
\thispagestyle{fancy}

\begin{abstract}
\noindent 這是摘要。本樣板使用 \textbf{TW-Kai(全字庫正楷體)} 作為中文字型,
並透過 \texttt{AutoFakeBold} 讓 \textbf{中文粗體}正常顯示。以下示範
\textbf{粗體}、\emph{斜體}、行內數學 $f_L=\gamma_p B/2\pi$ 與單位
\SI{2.675e8}{\radian\per\second\per\tesla}。
\end{abstract}

\section{前言}
這是一段內文。重點以\textbf{粗體}標示,例如\textbf{自旋--晶格弛豫時間 $T_1$}
與\textbf{表觀橫向弛豫時間 $T_2^{*}$}。可引用文獻~\cite{Sample2024}。

\section{方法與公式}
方程式範例(可用 \texttt{\textbackslash eqref} 交叉參照):
\begin{equation}
M(t)=M_\infty\!\left(1-e^{-t/T_1}\right).
\label{eq:sat}
\end{equation}
如式~\eqref{eq:sat} 所示,磁化量隨時間飽和。

\begin{thm}[範例定理]
這是定理環境的範例(標題為粗體)。
\end{thm}

\section{圖與表}
% --- 圖(實際使用時取消註解,並把圖檔放進 figures/)---
% \begin{figure}[htbp]
%   \centering
%   \includegraphics[width=0.7\linewidth]{your_figure.png}
%   \caption{圖說。}
%   \label{fig:demo}
% \end{figure}
\begin{figure}[htbp]
  \centering
  \fbox{\rule{0pt}{3cm}\rule{6cm}{0pt}}%
  \caption{圖片佔位框(把上面那行換成 \texttt{\textbackslash includegraphics}）。}
  \label{fig:demo}
\end{figure}

表格範例(\texttt{booktabs} 風格)見表~\ref{tab:demo}。
\begin{table}[htbp]
  \centering
  \caption{表格範例。}
  \label{tab:demo}
  \begin{tabular}{lcc}
    \toprule
    項目 & 量測值 & 不確定度 \\
    \midrule
    $T_1$ (s)     & $3.74$ & $0.18$ \\
    $T_2^{*}$ (s) & $0.75$ & ---    \\
    \bottomrule
  \end{tabular}
\end{table}

\section{結論}
\textbf{結論}:本樣板可獨立編譯,中文為 TW-Kai 且\textbf{粗體有效}。

\printbibliography
\end{document}
